Este trabalho comparou três modelos de tradução automática baseados na arquitetura Transformer, considerando a tradução do inglês para o português, a qualidade das saídas e a eficiência computacional. Nas condições experimentais estabelecidas, o MarianMT apresentou o melhor desempenho geral, enquanto os modelos M2M-100 e NLLB apresentaram resultados inferiores e mais próximos entre si.

Os cinco objetivos específicos foram atendidos. Primeiro, foram caracterizados o MarianMT, modelo bilíngue e especializado, e os modelos multilíngues M2M-100 e NLLB. Segundo, o MarianMT obteve os maiores valores de BLEU, BERTScore e chrF nos textos gerais e específicos. Terceiro, apresentou os menores valores de uso médio de CPU e tempo de inferência, enquanto os picos de memória foram próximos entre os modelos. Quarto, observou-se que os textos específicos produziram resultados inferiores aos textos gerais nas três métricas de qualidade. Quinto, verificou-se que BLEU, BERTScore e chrF indicaram o melhor desempenho do MarianMT, embora tenham apresentado diferenças na ordenação entre M2M-100 e NLLB.

Em resposta ao problema de pesquisa, o MarianMT apresentou o desempenho mais favorável no cenário analisado, pois combinou maior qualidade de tradução, menor uso médio de CPU e menor tempo de inferência. O pico de memória não diferenciou significativamente os modelos. Portanto, nas condições deste experimento, o MarianMT mostrou o melhor compromisso entre qualidade e eficiência computacional.

\subsection{Limitações}

As frases tinham aproximadamente seis palavras.

A execução experimental ocorreu somente em CPU.

Comparou-se um modelo bilíngue com modelos multilíngues.

Essas condições restringem a generalização dos resultados, pois não representam todos os tamanhos de sentença, ambientes computacionais ou configurações de modelos.

\subsection{Trabalhos futuros}

Algumas possibilidades de continuidade são apresentadas a seguir:

\begin{itemize}
\item Avaliar frases com diferentes comprimentos e domínios textuais;

\item Investigar outras arquiteturas de tradução automática;

\item Comparar modelos de pesos fechados e abertos;

\item Avaliar os modelos utilizando GPU e diferentes ambientes computacionais.
\end{itemize}
